% TEMPLATE-MILC-v3.tex - plantilla maestra para todas las guias MILC v3
% Ver scripts/generadores/build-guia-milc.py para la lista completa de
% claves esperadas. El script valida que no queden placeholders sin
% reemplazar antes de escribir el archivo final.

\documentclass[11pt,letterpaper]{article}
\usepackage[margin=1.75cm]{geometry}
\usepackage[table]{xcolor}
\usepackage{fontspec}
\usepackage[spanish]{babel}
\usepackage{tikz}
\usepackage[most]{tcolorbox}
\usepackage{tabularx}
\usepackage{array}
\usepackage{fancyhdr}
\usepackage{titlesec}
\usepackage{ragged2e}
\usepackage{enumitem}
\usepackage{microtype}

\setmainfont{Helvetica Neue}
\setsansfont{Helvetica Neue}
\setlength{\parindent}{0pt}
\setlength{\parskip}{6pt}
\hyphenpenalty=200
\exhyphenpenalty=200
\pretolerance=400
\tolerance=2000
\setlength{\emergencystretch}{1em}
\renewcommand{\arraystretch}{1.25}
\setlist{leftmargin=5mm,itemsep=1mm,topsep=1mm}
\newcolumntype{Y}{>{\RaggedRight\arraybackslash}X}

\definecolor{milcMagenta}{HTML}{E6007E}
\definecolor{milcVino}{HTML}{5A0038}
\definecolor{milcTurquesa}{HTML}{0093A5}
\definecolor{milcVerde}{HTML}{58A923}
\definecolor{milcMostaza}{HTML}{E5B400}
\definecolor{milcOcre}{HTML}{B86600}
\definecolor{milcNegro}{HTML}{241816}
\definecolor{milcGris}{HTML}{F7F7F4}
\definecolor{milcLinea}{HTML}{E8E2DD}
\definecolor{milcGradePrimary}{HTML}{FF2D87}
\definecolor{milcGradeDark}{HTML}{9F1B56}
\definecolor{milcGradeSoft}{HTML}{FFE0EE}
\definecolor{milcGradeAccent}{HTML}{CFE7FF}
\definecolor{milcGradeText}{HTML}{FFFFFF}
\color{milcNegro}

\pagestyle{fancy}
\fancyhf{}
\lhead{\small Tecnología e Informática · Grado 10}
\rhead{\small Guía 29 · MILC v3}
\cfoot{\small\thepage}
\renewcommand{\headrulewidth}{0pt}

\newtcolorbox{titlebox}[1]{
  enhanced,colback=#1,colframe=#1,arc=7mm,boxrule=0pt,
  left=7mm,right=7mm,top=3mm,bottom=3mm,
  before skip=8pt,after skip=8pt,
  fontupper=\sffamily\bfseries\Large\color{white}
}

\newtcolorbox{softbox}[3]{
  enhanced,colback=#2,colframe=#1,borderline west={4pt}{0pt}{#1},
  arc=2mm,boxrule=.4pt,left=4mm,right=4mm,top=3mm,bottom=3mm,
  before skip=6pt,after skip=8pt,title={#3},coltitle=white,
  colbacktitle=#1,fonttitle=\sffamily\bfseries,fontupper=\RaggedRight,
  attach boxed title to top left={xshift=3mm,yshift=-2mm},
  boxed title style={arc=2mm, boxrule=0pt}
}

\newtcolorbox{drawbox}[2]{
  enhanced,colback=white,colframe=#1,arc=4mm,boxrule=1pt,
  left=5mm,right=5mm,top=4mm,bottom=4mm,before skip=8pt,after skip=8pt,
  title={#2},coltitle=white,colbacktitle=#1,fonttitle=\sffamily\bfseries,
  fontupper=\RaggedRight,
  attach boxed title to top left={xshift=4mm,yshift=-2mm},
  boxed title style={arc=3mm, boxrule=0pt}
}

\newtcolorbox{infoband}[1]{
  enhanced,colback=#1!8,colframe=#1!50,arc=2mm,boxrule=.5pt,
  left=4mm,right=4mm,top=2mm,bottom=2mm,before skip=4pt,after skip=4pt,
  fontupper=\small\RaggedRight
}

% Puente narrativo entre secciones (contrato editorial Conéctate).
% Da continuidad: "Acabas de X. Ahora vas a Y porque Z."
% Visualmente: banda izquierda en color del grado, texto en itálica pequeña.
\newtcolorbox{puentebox}{
  enhanced,colback=milcGradeSoft,colframe=milcGradeDark,
  borderline west={3pt}{0pt}{milcGradeDark},
  arc=0mm,boxrule=0pt,left=5mm,right=4mm,top=2.5mm,bottom=2.5mm,
  before skip=4pt,after skip=8pt,
  fontupper=\itshape\small\color{milcNegro}\RaggedRight
}
\newcommand{\puente}[1]{%
  \begin{puentebox}%
    \textcolor{milcGradeDark}{\textbf{\textrightarrow}}\hspace{2mm}#1%
  \end{puentebox}%
}

\newcommand{\linead}[1]{\noindent\color{milcLinea}\rule{#1}{0.5pt}\color{milcNegro}}
\newcommand{\respuesta}[1]{\par\vspace{2mm}\foreach \n in {1,...,#1}{\noindent\color{milcLinea}\rule{\linewidth}{0.5pt}\par\vspace{5mm}}\color{milcNegro}}
\newcommand{\checkbox}{\textcolor{milcGradeDark}{\fbox{\phantom{x}}}\hspace{5pt}}

\newcommand{\phasecard}[4]{
\begin{tcolorbox}[enhanced,colback=#3,colframe=#2,arc=3mm,boxrule=.5pt,height=3.05cm,valign=center,halign=center,left=1.5mm,right=1.5mm,top=2mm,bottom=2mm]
{\sffamily\bfseries\fontsize{22}{24}\selectfont\textcolor{#2}{#1}}\par
{\sffamily\bfseries\small #4}
\end{tcolorbox}
}

\begin{document}
\thispagestyle{empty}
\begin{tikzpicture}[remember picture,overlay]
  \fill[white] (current page.south west) rectangle (current page.north east);
  \path[fill=milcGradePrimary,rounded corners=16mm]
    ([xshift=.55cm,yshift=-.55cm]current page.north west) rectangle
    ([xshift=-.55cm,yshift=3.65cm]current page.south east);

  \node[anchor=north west,text=milcGradeText] at ([xshift=2.0cm,yshift=-1.85cm]current page.north west)
    {\sffamily\bfseries\fontsize{14}{16}\selectfont GRADO};
  \node[anchor=north west,text=milcGradeText,opacity=.82] at ([xshift=2.0cm,yshift=-2.75cm]current page.north west)
    {\sffamily\bfseries\fontsize{9}{11}\selectfont SERIE GUÍAS · TECNOLOGÍA E INFORMÁTICA};

  \begin{scope}[shift={([xshift=-3.05cm,yshift=-2.55cm]current page.north east)}]
    \fill[white,opacity=.80] (-.62,-.52) rectangle (.62,.52);
    \draw[milcGradeText,line width=1pt] (-.62,-.52) rectangle (.62,.52);
    \fill[milcGradeDark] (-.42,-.38) rectangle (-.22,.06);
    \fill[milcGradeAccent] (-.10,-.38) rectangle (.10,.24);
    \fill[milcGradePrimary!60!black] (.22,-.38) rectangle (.42,.38);
  \end{scope}

  \node[anchor=west,text=milcGradeText] at ([xshift=1.9cm,yshift=-12.85cm]current page.north west)
    {\sffamily\bfseries\fontsize{124}{124}\selectfont 10};
  \draw[milcGradeText,line width=1.7pt]
    ([xshift=4.52cm,yshift=-11.68cm]current page.north west) circle (.13cm);

  \node[anchor=west,text=milcGradeText,text width=8.7cm,align=left] at ([xshift=10.35cm,yshift=-11.6cm]current page.north west)
    {
     {\sffamily\bfseries\fontsize{19}{22}\selectfont Proyecto integrador ---\\micro-emprendimiento\\con plan completo}\\[4mm]
     {\sffamily\bfseries\fontsize{23}{26}\selectfont Guía 29-10-TIC}\\[2mm]
     {\sffamily\fontsize{13}{16}\selectfont Período 3 · Ofimática con IA: contabilidad, datos y emprendimiento}\\[5mm]
     {\sffamily\bfseries\fontsize{11}{13}\selectfont Institución Educativa Sor María Juliana}\\[1mm]
     {\sffamily\fontsize{10}{12}\selectfont Autor: Dr. Álvaro Cárdenas Orozco}
    };

  \path[fill=milcGradeSoft,rounded corners=7mm]
    ([xshift=1.35cm,yshift=.85cm]current page.south west) rectangle
    ([xshift=-1.35cm,yshift=4.25cm]current page.south east);
  \draw[milcGradeDark,line width=.65pt,rounded corners=7mm]
    ([xshift=1.35cm,yshift=.85cm]current page.south west) rectangle
    ([xshift=-1.35cm,yshift=4.25cm]current page.south east);
  \node[anchor=center,text=milcNegro,text width=17.0cm,align=center] at ([yshift=2.55cm]current page.south)
    {\sffamily\fontsize{10}{12}\selectfont Metodología MILC · Apertura ancestral · Pensamiento computacional · Triángulo Dussel-Estoicismo-Floridi\\
    \textcolor{milcGradeDark}{\bfseries Producto:} Proyecto integrador final del grado 10° (puente directo al proyecto emprendedor del grado 11°): micro-emprendimiento personal o grupal completo en documento integrador de 8-12 páginas con (a) idea con problema real identificado y validado por estudio de mercado, (b) Business Model Canvas completo, (c) contabilidad proyectada a 3 meses con flujo de caja y punto de equilibrio, (d) las 3 piezas de comunicación empresarial (propuesta, correo, publicación), (e) plan de validación con 5 personas reales para próximas semanas, (f) bitácora del proceso completo del año};
\end{tikzpicture}
\mbox{}
\newpage

\section*{Apertura: Proyecto integrador --- micro-emprendimiento con plan completo}

\begin{softbox}{milcGradeDark}{milcGradeSoft}{Saber ancestral}
En los barrios populares del Valle del Cauca, en los corregimientos cafeteros del Quindío, en las cabeceras municipales del Pacífico colombiano, hay una figura que cualquier vecino mayor recuerda con respeto: \textbf{el aprendiz emprendedor que abre su primer puestico}. No el comerciante experimentado con 30 años de oficio; el joven o joven adulto que comienza. Ese aprendiz tenía una práctica que cualquier microempresario veterano del barrio reconocía: \textbf{no improvisaba}. Antes de poner el primer puestico de empanadas en la esquina, antes de abrir la primera papelería, antes de salir con la primera carga de aguacates a vender, el aprendiz \textbf{preparaba 5 cosas con disciplina}: (1) \textbf{Estudiaba la calle}: a qué hora pasa más gente, qué edad tienen, qué compran. (2) \textbf{Calculaba los costos}: cuánto cuesta cada empanada producir, cuánto la papel, cuánto el cajón. (3) \textbf{Ahorraba el capital inicial}: no se lanzaba con préstamo caro; ahorraba al menos lo necesario para 2 meses sin venta. (4) \textbf{Preparaba el lugar}: permiso del dueño, espacio limpio, presentación cuidada. (5) \textbf{Hablaba con vecinos}: les contaba la idea, escuchaba sus consejos. Esa preparación ancestral del aprendiz emprendedor es \textbf{exactamente lo que el proyecto integrador del grado 10° formaliza}. No improvisar es disciplina del oficio.
\end{softbox}

\begin{softbox}{milcTurquesa}{milcGris}{Saber contemporáneo}
El \textbf{proyecto integrador del periodo 3} es la pieza final del grado 10°: el aprendiz integra todas las habilidades del año (escritura con IA, informes técnicos, ofimática contable, estudio de mercado, Canvas, comunicación empresarial) en un \textbf{plan completo de micro-emprendimiento}. Sus 6 componentes obligatorios: (1) \textbf{Idea con problema real}: validada por estudio de mercado (S6) con datos verificables. (2) \textbf{Business Model Canvas completo} (S7) con los 9 bloques. (3) \textbf{Contabilidad proyectada a 3 meses} (S4) con flujo de caja, punto de equilibrio, 3 escenarios y decisión de viabilidad. (4) \textbf{Las 3 piezas de comunicación empresarial} (S8): propuesta, correo, publicación. (5) \textbf{Plan de validación}: 5 personas reales con quienes conversar en las próximas 2-3 semanas para verificar supuestos críticos del Canvas. (6) \textbf{Bitácora del proceso del año}: 1-2 páginas reflexionando sobre lo aprendido en los 3 periodos. La regla profesional: \textbf{este proyecto no es ejercicio escolar; es preparación seria para iniciar microempresa o como puente al proyecto emprendedor del grado 11°}.
\end{softbox}

\begin{softbox}{milcMostaza}{milcGris}{Pregunta-puente}
\textit{¿Qué sabía el aprendiz emprendedor del barrio al preparar 5 cosas con disciplina antes de abrir su primer puestico, que el novato olvida cuando lanza microempresa solo con \emph{``ya veremos''}? ¿Y por qué el proyecto integrador del grado 10° es realmente preparación para la vida real, no solo nota escolar?}
\end{softbox}

\begin{softbox}{milcVerde}{milcGris}{Saber y saber hacer}
Al terminar podrás: (1) \textbf{identificar} cómo integrar las 8 sesiones del periodo en un documento coherente; (2) \textbf{analizar} qué componentes ya tienes producidos y cuáles requieren refinamiento; (3) \textbf{crear} el documento integrador completo con los 6 componentes; (4) \textbf{evaluar} la viabilidad real del micro-emprendimiento y el plan de validación de próximas semanas.
\end{softbox}

\small
\begin{tabularx}{\textwidth}{>{\bfseries}p{3.0cm}Y}
\rowcolor{milcGradePrimary!12}Autor & Dr. Álvaro Cárdenas Orozco\\
DBA & Aplicación de ofimática asistida por IA a contabilidad básica y emprendimiento (MEN, T\&I)\\
Referentes & Lean Startup · Phronesis financiera · Ética del dato empresarial\\
Producto final & Proyecto integrador final del grado 10° (puente directo al proyecto emprendedor del grado 11°): micro-emprendimiento personal o grupal completo en documento integrador de 8-12 páginas con (a) idea con problema real identificado y validado por estudio de mercado, (b) Business Model Canvas completo, (c) contabilidad proyectada a 3 meses con flujo de caja y punto de equilibrio, (d) las 3 piezas de comunicación empresarial (propuesta, correo, publicación), (e) plan de validación con 5 personas reales para próximas semanas, (f) bitácora del proceso completo del año\\
\end{tabularx}
\normalsize

\puente{Antes de empezar, mira el viaje completo. Has recorrido 8 sesiones del periodo 3. Hoy es el día de \textbf{integrar el año entero} en una sola pieza maestra. La sesión 10 será sustentación pública + feria de microempresas, cierre del grado 10°.}

\begin{titlebox}{milcGradeDark}
Ruta MILC: así vamos a aprender
\end{titlebox}
\begin{tabularx}{\textwidth}{>{\centering\arraybackslash}X>{\centering\arraybackslash}X>{\centering\arraybackslash}X>{\centering\arraybackslash}X}
\phasecard{1}{milcVerde}{milcGris}{Escuta\\[-1mm]\small Observo, escucho y pregunto.} &
\phasecard{2}{milcTurquesa}{milcGris}{Sistematización\\[-1mm]\small Pensamiento computacional.} &
\phasecard{3}{milcMagenta}{milcGris}{Praxis\\[-1mm]\small Construyo y aplico.} &
\phasecard{4}{milcVino}{milcGris}{Evaluación\\[-1mm]\small Reflexión liberadora.}
\end{tabularx}


\puente{Antes de teorizar, vas a hacer un \emph{inventario completo}: lista lo que tienes producido del periodo 3 (cuaderno de cuentas, plantilla Sheets, proyección de flujo, gráficos, estudio mercado, Canvas, 3 piezas comunicación) y lo que falta para el documento integrador.}

\begin{titlebox}{milcVerde}
Fase 1 · Escuta
\end{titlebox}

\begin{softbox}{milcVerde}{milcGris}{Escena para escuchar}
\textbf{Actividad 1 · IDENTIFICA --- Inventario completo del periodo 3} (15 min · individual). Haz lista de lo que tienes producido en las 8 sesiones del periodo 3. Reglas: (1) \textbf{Tengo}: mapa ofimática (S1), cuaderno cuentas 2 semanas (S2), plantilla Sheets (S3), proyección flujo 3 meses (S4), 3 gráficos honestos (S5), estudio mercado con hallazgos (S6), Canvas completo (S7), 3 piezas comunicación (S8). (2) \textbf{Me falta}: integración en documento único, plan de validación de 5 personas, bitácora del año.
\end{softbox}

\checkbox Listé lo producido en las 8 sesiones.\\
\checkbox Listé lo que falta para integrar.\\
\checkbox Estimé tiempo necesario para cerrar.

\begin{infoband}{milcVerde}
\textbf{Lo que va al cuaderno (Actividad 1 · IDENTIFICA).} Título: «Actividad 1 --- Inventario completo». \textbf{Formato:} 2 listas con cronograma. \textbf{Sabes que terminaste cuando} tienes plan claro de integración.
\end{infoband}

\puente{Ya tienes inventario. Ahora vas a entender la \textbf{anatomía del proyecto integrador} y la cadena lógica entre componentes.}

\begin{titlebox}{milcTurquesa}
Fase 2 · Sistematización con pensamiento computacional
\end{titlebox}

\begin{softbox}{milcTurquesa}{milcGris}{Cuatro pilares aplicados al tema}
Tu inventario es el plan. Ahora necesitas la \textbf{anatomía del proyecto integrador}.
\end{softbox}

\small
\begin{tabularx}{\textwidth}{>{\bfseries\RaggedRight\arraybackslash}p{3.6cm}Y}
\rowcolor{milcTurquesa!90}\color{white}Pilar & \color{white}\bfseries Aplicación al tema\\
1. Descomposición & Los \textbf{6 componentes obligatorios} se descomponen así: (1) \textbf{Idea con problema real}: 1 página presentando el microemprendimiento, problema validado, audiencia. (2) \textbf{Business Model Canvas}: imagen o tabla del Canvas con los 9 bloques. (3) \textbf{Contabilidad proyectada}: tabla de flujo de caja 3 meses, punto de equilibrio, decisión de viabilidad. (4) \textbf{Comunicación empresarial}: las 3 piezas (propuesta comercial, correo, publicación visual). (5) \textbf{Plan de validación}: lista de 5 personas reales con nombres y contexto para conversar en las próximas 2-3 semanas verificando los supuestos críticos del Canvas. (6) \textbf{Bitácora del proceso del año}: 1-2 páginas reflexivas sobre lo aprendido en los 3 periodos del grado 10° (libro con IA, informes profesionales, ofimática emprendedora).\\
2. Reconocimiento de patrones & La \textbf{cadena lógica} entre componentes es crucial: (a) Idea sostenida por estudio mercado (S6). (b) Canvas alimentado por estudio mercado en al menos 3 bloques. (c) Contabilidad coherente con Canvas (precios, costos, escala). (d) Comunicación que refleja la propuesta de valor del Canvas. (e) Plan de validación dirigido a supuestos críticos del Canvas. Si rompes la cadena, el proyecto parece pero no es coherente.\\
3. Abstracción & Los \textbf{5 errores típicos} del proyecto integrador novato: (a) \textbf{Piezas yuxtapuestas sin integración}: pegar Canvas + flujo + comunicación sin conectarlos. (b) \textbf{Plan de validación con familia solamente}: incluir solo amigos cercanos. (c) \textbf{Sin bitácora del año}: el documento llega sin reflexión sobre lo aprendido. (d) \textbf{Decisión de viabilidad inflada}: pretender viabilidad cuando el flujo dice lo contrario. (e) \textbf{Sin uso explícito de IA documentado}: ocultar que hubo asistencia de IA.\\
4. Algoritmo conceptual & Algoritmo: (1) crear documento integrador (Google Docs profesional de S3 del P2). (2) sección 1: idea + problema. (3) sección 2: Canvas. (4) sección 3: contabilidad proyectada. (5) sección 4: 3 piezas de comunicación. (6) sección 5: plan de validación 5 personas. (7) sección 6: bitácora del año. (8) revisión final de coherencia.\\
\end{tabularx}
\normalsize

\begin{drawbox}{milcTurquesa}{Anatomía del proyecto integrador}
\textbf{6 componentes:} idea / Canvas / contabilidad / comunicación / validación / bitácora. \\[1mm] \textbf{Cadena lógica:} cada pieza sostiene la siguiente. \\[1mm] \textbf{Extensión:} 8-12 páginas. \\[1mm] \textbf{Plan validación:} 5 personas reales con nombres. \\[1mm] \textbf{Bitácora del año:} reflexión personal.
\end{drawbox}

\begin{softbox}{milcMostaza}{milcGris}{Errores comunes y su corrección}
(1) Piezas yuxtapuestas sin integración. (2) Validación solo con familia cercana. (3) Sin bitácora del año. (4) Decisión de viabilidad inflada. (5) Sin uso explícito de IA documentado. (6) Sin plan concreto para próximas semanas.
\end{softbox}

\begin{infoband}{milcTurquesa}
\textbf{Lo que va al cuaderno (Actividad 2 · ANALIZA).} Título: «Actividad 2 --- Anatomía proyecto integrador». \textbf{Formato:} (a) 6 componentes con función; (b) cadena lógica; (c) 6 errores con marca. \textbf{Sabes que terminaste cuando} sabes integrar el año en 1 documento coherente.
\end{infoband}

\puente{Tienes el método. Toca aplicarlo: integras componentes en documento único, agregas plan de validación de 5 personas, escribes bitácora del año.}

\begin{titlebox}{milcMagenta}
Fase 3 · Praxis con pensamiento computacional
\end{titlebox}

\begin{softbox}{milcMagenta}{milcGris}{Construyendo el producto con los cuatro pilares}
\textbf{Actividad 3 · CREA + EVALÚA --- Documento integrador completo} (60 min en sesión + tiempo en casa · individual o parejas). Hora de integrar. Construyes documento único de 8-12 páginas con los 6 componentes.
\end{softbox}

\small
\begin{tabularx}{\textwidth}{>{\bfseries\RaggedRight\arraybackslash}p{3.6cm}Y}
\rowcolor{milcMagenta!90}\color{white}Pilar & \color{white}\bfseries Aplicación al producto\\
Descomposición & Tu documento se construye en 8 pasos: (1) abrir Google Docs con plantilla profesional. (2) escribir Sección 1: Idea + problema validado (1-2 páginas). (3) integrar Sección 2: Canvas (imagen del Canvas + explicación de 1 página). (4) Sección 3: Contabilidad proyectada (2 páginas con flujo, PE, escenarios, decisión). (5) Sección 4: 3 piezas de comunicación (2-3 páginas). (6) Sección 5: Plan de validación con 5 personas reales (1 página). (7) Sección 6: Bitácora del año (1-2 páginas). (8) revisión de coherencia + exportar PDF.\\
Patrones & El \textbf{plan de validación con 5 personas reales} es la pieza puente al grado 11°. Tiene 5 filas (1 por persona) con: nombre, contexto (¿es cliente potencial?, ¿es socio potencial?, ¿es experto?), qué supuesto crítico del Canvas verificarás con esa persona, cómo contactarás, cuándo conversarás. La phronesis del oficio consiste en \textbf{tener nombres reales contactables, no \emph{``algún emprendedor que conozca''}}.\\
Abstracción & La \textbf{bitácora del año} responde 5 preguntas reflexivas: (a) ¿qué aprendí del P1 (escritura con IA)?. (b) ¿qué del P2 (informes profesionales)?. (c) ¿qué del P3 (ofimática y emprendimiento)?. (d) ¿qué cambió en mi forma de usar IA durante el año?. (e) ¿qué me llevo al grado 11°? Esa bitácora es la firma personal del año.\\
Algoritmo de armado & Algoritmo: (1) abrir Docs. (2) las 6 secciones en orden. (3) plan validación con nombres. (4) bitácora del año. (5) revisión coherencia. (6) PDF exportado.\\
\end{tabularx}
\normalsize

\begin{drawbox}{milcMagenta}{Checklist al cerrar el proyecto}
\checkbox Idea con problema validado por estudio mercado. \\[1mm] \checkbox Canvas completo integrado. \\[1mm] \checkbox Contabilidad proyectada 3 meses con decisión. \\[1mm] \checkbox 3 piezas comunicación coherentes. \\[1mm] \checkbox Plan validación con 5 nombres reales. \\[1mm] \checkbox Bitácora del año completa. \\[1mm] \checkbox 8-12 páginas en PDF exportado.
\end{drawbox}

\begin{softbox}{milcMagenta}{milcGris}{Plantilla de trabajo}
\textbf{1. Idea + problema.} \textit{1-2 páginas.} \\[1mm] \textbf{2. Canvas.} \textit{Imagen + 1 página.} \\[1mm] \textbf{3. Contabilidad.} \textit{Flujo + PE + decisión, 2 páginas.} \\[1mm] \textbf{4. Comunicación.} \textit{3 piezas, 2-3 páginas.} \\[1mm] \textbf{5. Plan validación.} \textit{5 personas reales.} \\[1mm] \textbf{6. Bitácora del año.} \textit{Reflexión personal.}
\end{softbox}

\begin{infoband}{milcMagenta}
\textbf{Lo que va al cuaderno (Actividad 3 · CREA + EVALÚA).} Título: «Actividad 3 --- Proyecto integrador G10». \textbf{Formato:} (a) link al PDF; (b) verificación de los 6 componentes; (c) plan validación con nombres. \textbf{Sabes que terminaste cuando} el documento sirve realmente para iniciar microempresa.
\end{infoband}

\puente{Lo que sigue resume \emph{qué} entregas hoy: documento integrador 8-12 páginas + plan validación + bitácora del año. El criterio principal: que el documento sirva realmente como base para iniciar microempresa o como puente al proyecto emprendedor del grado 11°.}

\begin{titlebox}{milcMostaza}
Producto de la sesión
\end{titlebox}
\textbf{Proyecto integrador 8-12 pp + plan validación + bitácora del año}.
\begin{softbox}{milcMostaza}{milcGris}{Criterios de calidad}
(1) 6 componentes obligatorios presentes. (2) Cadena lógica entre piezas verificable. (3) Plan validación con 5 nombres reales. (4) Bitácora del año completa. (5) Coherencia entre Canvas, contabilidad y comunicación. (6) PDF exportado profesional.
\end{softbox}

\puente{Antes de cerrar, mira el proyecto integrador desde las cinco dimensiones humanas. Integrar con disciplina es respeto por el oficio del año entero; entregar piezas sueltas es desperdicio de 8 sesiones de trabajo.}

\begin{titlebox}{milcVino}
Fase 4 · Evaluación liberadora — cinco dimensiones
\end{titlebox}

\begin{softbox}{milcVino}{milcGris}{Sentido de la evaluación}
Evaluar no es solo revisar una nota. Es reconocer qué aprendí, qué cambió en mí, qué emociones manejé, cómo me relaciono con la ciudad y la comunidad, y qué le dejaré a quien viene detrás.
\end{softbox}

\textbf{Mi reflexión liberadora en cinco dimensiones.} Completa cada frase iniciada:

\small
\begin{tabularx}{\textwidth}{>{\bfseries\RaggedRight\arraybackslash}p{3.4cm}Y>{\RaggedRight\arraybackslash}p{4.6cm}}
\rowcolor{milcVino}\color{white}Dimensión & \color{white}\bfseries Mi respuesta (frame starter) & \color{white}\bfseries Algo que descubrí\\
1. Personal & Lo que más cambió en mí fue \linead{6.5cm} & \linead{4.2cm}\\
2. Emocional & La emoción más fuerte fue \linead{2.5cm} y la regulé \linead{3cm} & \linead{4.2cm}\\
3. Ciudadana & En democracia esto importa porque \linead{6cm} & \linead{4.2cm}\\
4. Local (Cartago) & En mi barrio sería útil para \linead{6.5cm} & \linead{4.2cm}\\
5. Intergeneracional & A mi abuela/abuelo le diría \linead{6.5cm} & \linead{4.2cm}\\
\end{tabularx}
\normalsize

\begin{infoband}{milcVino}
\textbf{Para la próxima sustentación.} Vuelve a esta tabla en seis meses. Verás que las respuestas cambian — no porque lo aprendido cambie, sino porque tú cambias. La evaluación liberadora es una conversación contigo a través del tiempo.
\end{infoband}


\puente{Y para terminar, tres voces sobre la pieza maestra emprendedora. Dussel sobre los micro-emprendimientos que dan voz a comunidades, Marco Aurelio sobre la integración con disciplina, Floridi sobre el proyecto como acto profesional contemporáneo.}

\begin{titlebox}{milcGradeDark}
Triángulo de pensamiento · cierre filosófico
\end{titlebox}

\begin{softbox}{milcVino}{milcGris}{Dussel · lente del nosotros}
\textit{``Un proyecto emprendedor que sirve a la comunidad y se sostiene financieramente es liberador; uno que solo extrae es opresor.''}

Tu micro-emprendimiento puede ser una de 2 cosas: (a) servicio real al barrio con sostenibilidad económica, (b) extracción disfrazada de emprendimiento. La filosofía de la liberación pide la primera. La phronesis del oficio consiste en \textbf{diseñar emprendimiento que dignifique al cliente y sostenga al emprendedor}.

\textbf{¿Mi micro-emprendimiento sirve realmente o solo busca extraer valor?}
\end{softbox}
\linead{\linewidth}\\[1mm]
\linead{\linewidth}

\begin{softbox}{milcMostaza}{milcGris}{Estoicismo · Marco Aurelio · cuidado interior}
\textit{``Integrar el año con disciplina es virtud emprendedora; pegar piezas sin coherencia es vanidad.''}

Es tentador pegar los 6 componentes en orden sin integrar. La sabiduría estoica recomienda lo opuesto: \emph{teje las piezas en un solo argumento coherente}. La phronesis del oficio honra la integración sobre la agregación.

\textbf{¿Mi proyecto integrador es realmente integrado, o solo es agregación de 6 piezas?}
\end{softbox}
\linead{\linewidth}\\[1mm]
\linead{\linewidth}

\begin{softbox}{milcTurquesa}{milcGris}{Floridi · lente de la infoesfera}
\textit{``El proyecto emprendedor con IA documentada honestamente es el modelo profesional emergente.''}

En la era de la IA generativa, los micro-emprendimientos que declaran abiertamente el uso de IA se posicionan como ético contemporáneo. Tu proyecto te entrena en una práctica que rige todo emprendimiento serio: \textbf{transparencia con porcentajes y modelos}.

\textbf{¿Mi proyecto contribuye al modelo emprendedor transparente del siglo XXI?}
\end{softbox}
\linead{\linewidth}\\[1mm]
\linead{\linewidth}

\begin{titlebox}{milcVino}
Compromiso final
\end{titlebox}

\small
\begin{tabularx}{\textwidth}{>{\RaggedRight\arraybackslash}p{3.0cm}YYY}
\rowcolor{milcVino}\color{white}\bfseries Criterio & \color{white}\bfseries Lo logré & \color{white}\bfseries En proceso & \color{white}\bfseries Necesito apoyo\\
Comprensión del tema & Explico ideas centrales con mis palabras. & Reconozco ideas pero falta precisión. & Necesito repasar conceptos base.\\
Pensamiento computacional & Apliqué los 4 pilares al diseñar mi producto. & Apliqué algunos pilares. & Necesito releer la fase 2.\\
Aplicación y producto & Mi producto cumple los criterios y se puede revisar. & Producto funcional con ajustes pendientes. & Aún en construcción.\\
Reflexión 5 dimensiones & Respondí cada dimensión con honestidad. & Respondí algunas con detalle. & Mis respuestas son muy generales.\\
Triángulo de pensamiento & Conecté las 3 voces con mi trabajo real. & Conecté al menos una. & No relaciono las voces con mi proceso.\\
\end{tabularx}
\normalsize

\textbf{Hoy entendí que:}\\[1mm]
\linead{\linewidth}\\[1mm]
\linead{\linewidth}

\textbf{Mi compromiso para la próxima sesión es:}\\[1mm]
\linead{\linewidth}\\[1mm]
\linead{\linewidth}

\begin{softbox}{milcMostaza}{milcGris}{Cita de despedida · Marco Aurelio}
\textit{``El obstáculo en el camino se vuelve el camino. Lo que estorba al camino se vuelve camino.''}\\[1mm]
Cada error de hoy, bien observado, se convierte en pieza del camino que pisarás mañana.
\end{softbox}

\small
\begin{tabularx}{\textwidth}{>{\bfseries}p{3.6cm}Y}
\rowcolor{milcGradeDark!12}Nombre completo: & \linead{10cm}\\
Fecha: & \linead{4cm} \quad Curso/grupo: \linead{4cm}\\
Mi firma: & \linead{9cm}\\
\end{tabularx}
\normalsize

\begin{infoband}{milcGradeDark}
\textbf{Para la próxima sesión.} Llega con tu proyecto integrador completo en PDF, plan de validación con 5 nombres y bitácora del año. En la sesión 10 vas a hacer la \textbf{sustentación pública del grado 10° + feria de microempresas escolares}: cierre del grado y puente al proyecto emprendedor del grado 11°.
\end{infoband}

\end{document}
